\section{Additional Settings and Ablations}
\label{app:settings}

The main paper reports the settings that the baselines' papers use: CogVideoX-2B at 49 frames and Wan2.1 at 81 frames. This appendix reports the third frame count we ran, so that the ordering can be checked away from the flagship configurations, and the ablations that fix the two parameters of the schedule.

\subsection{17 frames, 100 prompts}
\label{app:cog17}

At 17 frames the trajectory is cheaper per call, which lets us run all four methods over the full prompt set.

\begin{table}[h]
\centering\footnotesize\setlength{\tabcolsep}{4pt}
\begin{tabular}{lcccc}
\toprule
Method & Calls & LPIPS $\downarrow$ & SSIM $\uparrow$ & PSNR $\uparrow$ \\
\midrule
TaylorSeer & 19 & 0.492 & 0.526 & 13.9 \\
TeaCache & 17 & 0.279 & 0.752 & 21.1 \\
RA-CFGCache & 19 & 0.159 & 0.895 & 27.7 \\
\method{} & 20 & \textbf{0.145} & \textbf{0.906} & \textbf{28.6} \\
\bottomrule
\end{tabular}
\caption{CogVideoX-2B, 17 frames, 50 steps, 100 prompts, one pipeline and one reference.}
\label{tab:cog17}
\end{table}

The ordering of Table~\ref{tab:cog17} is the same as in the main table, and the gap to the strongest baseline is of the same size, which is what a placement effect should look like: it does not depend on the frame count through any mechanism of its own, only through the profile, which we re-measure for each setting.

\subsection{Wan2.1: the remaining baselines}
\label{app:wanfull}

\begin{table}[h]
\centering\footnotesize\setlength{\tabcolsep}{4pt}
\begin{tabular}{lcccc}
\toprule
Method & Calls & LPIPS $\downarrow$ & SSIM $\uparrow$ & PSNR $\uparrow$ \\
\midrule
TaylorSeer & 19 & 0.379 & 0.620 & 17.6 \\
TeaCache & 19 & 0.199 & 0.827 & 24.0 \\
RA-CFGCache & 19 & 0.154 & 0.870 & 27.8 \\
\method{} & 21 & \textbf{0.134} & \textbf{0.888} & \textbf{28.9} \\
\bottomrule
\end{tabular}
\caption{Wan2.1-T2V-1.3B, 81 frames, 50 steps, 16 prompts. We report these at the prompt count each baseline needed to finish; the two-row comparison in the main table uses 65 prompts for the risk-based cache and our schedule.}
\label{tab:wanfull}
\end{table}

\subsection{The shape of the tail}
\label{app:shape}

At a fixed head and a fixed budget, the only freedom left is the spacing of the tail. With the head fixed at twelve steps and the budget at twenty calls, we vary only the ratio at which the gaps grow. On CogVideoX-2B at 17 frames over 32 prompts the uniform tail reaches 0.166 LPIPS, and the growing tails reach 0.156, 0.157 and 0.154 at ratios of 1.4, 2.0 and 3.0. The growth is therefore worth 5 to 7 percent, and its exact value is not: the three ratios are within noise of one another. The Taylor expansion baseline is swept over its own interval in the next subsection; that comparison is not evidence for the shape, because its degradation is dominated by its forecast.

\subsection{Why the forecast-based baseline degrades}
\label{app:taylor}

A reader may suspect that the Taylor expansion baseline is simply given a smaller budget than ours, and that a denser spacing would rescue it. Table~\ref{tab:taylor} rules this out: at an interval of two steps it spends about 26 calls, 30\% more than our twenty, and still sits at 0.402 LPIPS, almost three times our 0.142.

\begin{table}[h]
\centering\footnotesize\setlength{\tabcolsep}{5pt}
\begin{tabular}{lccc}
\toprule
Spacing & Calls & LPIPS $\downarrow$ & PSNR $\uparrow$ \\
\midrule
TaylorSeer, interval 2 & 26 & 0.402 & 16.7 \\
TaylorSeer, interval 3 & 19 & 0.531 & not measured \\
TaylorSeer, interval 4 & 15 & 0.615 & 11.6 \\
\midrule
\method{} & 20 & \textbf{0.142} & \textbf{30.2} \\
\bottomrule
\end{tabular}
\caption{The Taylor expansion baseline at three spacings, against our schedule, on CogVideoX-2B at 17 frames and 32 prompts. Denser spacing does not close the gap, so the loss is in the forecast rather than in the budget.}
\label{tab:taylor}
\end{table}

The mechanism is a mismatch between the extrapolator's assumption and the sampler it is applied to. A first-order Taylor expansion assumes the hidden state evolves smoothly between the steps it observes. CogVideoX-2B predicts $v$ rather than $\epsilon$, and with fifty steps its latent moves in jumps rather than along a smooth curve, so a derivative estimated from the two previous outputs points away from the next one. The released defaults of that baseline (an interval of three steps, first order, three enhanced steps) were tuned on architectures whose features are smoother, and its own paper does not report CogVideoX. This is why we treat it as a baseline rather than as a competing design point: its degradation is a property of the forecast, and a better forecast would be a different method, not a better schedule.

\subsection{The head length, measured}
\label{app:head}

Table~\ref{tab:head} varies the head at a fixed budget of about twenty calls, with the tail grown from the head to the last step in each case.

\begin{table}[h]
\centering\footnotesize\setlength{\tabcolsep}{5pt}
\begin{tabular}{lcccc}
\toprule
Head $H$ & Calls & LPIPS $\downarrow$ & SSIM $\uparrow$ & PSNR $\uparrow$ \\
\midrule
6  & 20 & 0.288 & 0.772 & 21.5 \\
12 & 20 & 0.157 & 0.907 & 28.3 \\
16 & 22 & \textbf{0.133} & \textbf{0.930} & \textbf{30.8} \\
\bottomrule
\end{tabular}
\caption{CogVideoX-2B, 17 frames, 32 prompts. A short head is catastrophic and a longer one keeps helping at this budget. At 49 frames and 100 prompts the same comparison gives 0.142 with a head of twelve and 0.130 with a head of sixteen, so the schedule we ship uses sixteen.}
\label{tab:head}
\end{table}

Two things follow. A head of six is far worse than one of twelve, which is the sharp end of the effect and the reason a uniform grid loses: it effectively takes a head of zero. And the optimum at a budget of twenty calls lies at or above sixteen, not at twelve, which means the head should be read from the budget rather than fixed. At 49 frames and 100 prompts this is confirmed end to end: a head of twelve reaches 0.142 LPIPS and a head of sixteen 0.130, at 20 and 22 calls respectively.

\subsection{The two parameters of the schedule}
\label{app:params}

The schedule has two parameters: the head length $H$ and the ratio at which the gaps grow. $H$ is not searched; it is read from the measured profile as the number of leading steps above a fraction of the maximum. Read at one half of the maximum, the profile gives four steps on CogVideoX-2B and more when the fraction is lowered; at one tenth it gives nine. The schedules we ship use twelve, where the measured profile has decayed to about a tenth of its maximum. The ordering against the baselines is preserved for every head we tried, so the exact value is not a knife edge. What the budget decides is how much of the head can be held exact.
